\chapter*{Preface}
\addcontentsline{toc}{chapter}{Preface}
\markboth{Preface}{}

In 1954 Chen Ning Yang and Robert Mills wrote down a generalisation of Maxwell's equations in which the field carries its own charge. The equations were beautiful and, for a decade, apparently useless: a classical Yang Mills field has massless waves, yet the strong force acts only over a range of about a fermi. The resolution, assembled over the next twenty years, is one of the great stories of physics. Quantised and placed on a lattice, the theory confines; at short distances it is asymptotically free; and at long distances its lightest excitation, a glueball\index{glueball}, has positive mass. That last statement, the mass gap\index{mass gap}, is the half of the Yang Mills Millennium Prize Problem that remains open.

The other half was reported to have changed status while this book was being written. On 26 October 2026 OpenAI announced~\cite{OAI} that one of its models had produced a proof of the existence of pure quantum Yang Mills theory in four dimensions. On 10 November 2026 a working group chaired by David Gross, with Michael Douglas, Arthur Jaffe, Igor Klebanov and Nikita Nekrasov among its members, released the Lean formalisation together with a validation assessment. The working group was explicit that confinement\index{confinement} and the mass gap remain open. This book is about the half that remains.

We come to the problem as statisticians. Our main research programme over the last two years has been a probabilistic attack on the Riemann Hypothesis\index{Riemann Hypothesis} through generalized gamma convolution\index{generalized gamma convolution}s, Thorin measure\index{Thorin measure}s and the transport of L\'evy measures into Thorin measures by what we call the $2\sin^2$ lemma. In that programme everything turns on positivity: the Riemann Hypothesis becomes the statement that a certain transported Thorin measure is nonnegative, and that is what we cannot yet prove. When we asked whether the same machinery says anything about Yang Mills, the answer surprised us. In Yang Mills, positivity is free. Reflection positivity, proved for lattice gauge theory by Osterwalder and Seiler~\cite{OSe}\index{Osterwalder, K.!and Seiler, E.} in 1978, makes every gauge invariant Euclidean autocorrelator, the correlator of an observable with its reflection, the Laplace transform of a nonnegative spectral measure\index{spectral measure}. The hard question moves from positivity to support: where does the measure live, and does its support stay away from zero as the lattice spacing goes to zero?

That role reversal organises the book. After stating the problem and the probabilistic tools, we give two exact dictionaries between the gap and Thorin measures: the spectral measure is a Thorin measure whose left edge is the gap, and the $2\sin^2$ lemma turns the spectrum into the jumps of a compound Poisson process whose L\'evy exponent is, up to the factor $\pi t$, a Thorin density decaying at the rate of the gap. In two dimensions, where the plaquette weight is the law of Brownian motion on SU(2), we prove a uniform gap inequality from the coarsest lattice to the continuum, and find Riemann's kernel $\Phi$ inside the return density of Brownian motion on the three sphere. In four dimensions the Migdal Kadanoff recursion becomes an alternation of Bayesian products and convolutions of laws on the group; flat space would leave the coupling marginal, and the curvature factor $\theta/\sin\theta$, which at imaginary angle is the Laplace transform of the Biane Pitman Yor variable $S_{\mathrm{BPY}}$, makes it run, to leading order, at rate $(1-b^{-2})N^2\E S_{\mathrm{BPY}}/4$ for SU(N), within three percent of one loop at $b=2$. We lay out a five stage path to the gap and the one comparison, between the Migdal Kadanoff ladder and the true lattice, on which it turns.

The later chapters work on that comparison. A product formula shows that the SU(2) heat kernel, as a function of $\cos\theta$, has only real zeros that move outward under the heat flow, which makes the heat kernel projected ladder monotone at every coupling. Tests against lattice data show what the ladder gets right, the rate of running to within fifteen percent, and what it gets wrong, the second order coefficient and the place of the crossover. And simulations in three and four dimensions show the true law of a blocked plaquette lying between two caricatures, independent plaquettes and Migdal Kadanoff, in the likelihood ratio order. That empirical sandwich is the comparison problem in its smallest form.

Along the way we ask of each of a series of ideas whether it is useful here: latent variables and nested sampling, the vertical likelihood, Yakowitz, Kac and P\'olya, Jerrum and Sinclair, Sokal, Besag, slice variables, and our own books on the Riemann Hypothesis and on Hilbert's sixth problem. The answer is yes more often than we expected.

We owe a particular debt to Mark Kac, who died in 1984. His commentary on P\'olya's paper on the $\xi$ function, which records how a lemma about the Riemann Hypothesis became the first proof of a case of the Lee and Yang theorem, is the hinge between our two programmes, and his method of the smallest solvable case, his formula with Feynman, his lemma on return times, his propagation of chaos and his essay on models are used throughout. Chapter~\ref{ch:kacyang} sets out what is his.

Theorems are proved; leading order derivations, numerical results and open problems are labelled as such. Every figure comes from the code in the appendix.

\vspace{1em}
\noindent Nick Polson and Vadim Sokolov\\
Chicago, November 2026
